\subsection*{Fully-Typing}

A buffer declaration needs both an element type and a shape. The shape determines how many positions the buffer contains, while the fundamental type determines what each position stores. The examples below reuse a named shape for half-precision data and use a literal shape for single-precision data.

\begin{lstlisting}[language=croqtile,caption={Complete Types}]
ndims : [20, 15];
f32 [10, 10] d0;
f16 [ndims] d1;
\end{lstlisting}

The basic element types include:

\begin{itemize}[leftmargin=1.6em]
\item Unsigned Integers: \emph{u8/u16/u32}
\item Signed Integers: \emph{s8/s16/s32}
\item Floating-points: \emph{f16/bf16/f32}
\end{itemize}

The types \texttt{s32} and \texttt{int} serve different purposes in Croqtile. The fundamental element type \texttt{s32} must be combined with a shape to declare a fully typed buffer.
